\documentclass[10pt,a4paper]{amsart}
\usepackage{amsmath,amssymb,amsthm,graphicx}
\usepackage{verbatim}
\newenvironment{reference}{\comment}{\endcomment}
%\newenvironment{reference}{}{}
\newenvironment{slogan}{\comment}{\endcomment}
\newenvironment{history}{\comment}{\endcomment}

% For commutative diagrams we use Xy-pic
\usepackage[all]{xy}

% We use 2cell for 2-commutative diagrams.
\xyoption{2cell}
% Theorem environments.
%
\theoremstyle{plain}
\newtheorem{theorem}[subsection]{Theorem}
\newtheorem{proposition}[subsection]{Proposition}
\newtheorem{lemma}[subsection]{Lemma}

\theoremstyle{definition}
\newtheorem{definition}[subsection]{Definition}
\newtheorem{example}[subsection]{Example}
\newtheorem{exercise}[subsection]{Exercise}
\newtheorem{situation}[subsection]{Situation}

\theoremstyle{remark}
\newtheorem{remark}[subsection]{Remark}
\newtheorem{remarks}[subsection]{Remarks}

\numberwithin{equation}{subsection}

% Macros
%
\def\lim{\mathop{\mathrm{lim}}\nolimits}
\def\colim{\mathop{\mathrm{colim}}\nolimits}
\def\Spec{\mathop{\mathrm{Spec}}}
\def\Hom{\mathop{\mathrm{Hom}}\nolimits}
\def\Ext{\mathop{\mathrm{Ext}}\nolimits}
\def\SheafHom{\mathop{\mathcal{H}\!\mathit{om}}\nolimits}
\def\SheafExt{\mathop{\mathcal{E}\!\mathit{xt}}\nolimits}
\def\Sch{\mathit{Sch}}
\def\Mor{\mathop{\mathrm{Mor}}\nolimits}
\def\Ob{\mathop{\mathrm{Ob}}\nolimits}
\def\Sh{\mathop{\mathit{Sh}}\nolimits}
\def\NL{\mathop{N\!L}\nolimits}
\def\CH{\mathop{\mathrm{CH}}\nolimits}
\def\proetale{{pro\text{-}\acute{e}tale}}
\def\etale{{\acute{e}tale}}
\def\QCoh{\mathit{QCoh}}
\def\Ker{\mathop{\mathrm{Ker}}}
\def\Im{\mathop{\mathrm{Im}}}
\def\Coker{\mathop{\mathrm{Coker}}}
\def\Coim{\mathop{\mathrm{Coim}}}

% Boxtimes
%
\DeclareMathSymbol{\boxtimes}{\mathbin}{AMSa}{"02}

%
% Macros for moduli stacks/spaces
%
\def\QCohstack{\mathcal{QC}\!\mathit{oh}}
\def\Cohstack{\mathcal{C}\!\mathit{oh}}
\def\Spacesstack{\mathcal{S}\!\mathit{paces}}
\def\Quotfunctor{\mathrm{Quot}}
\def\Hilbfunctor{\mathrm{Hilb}}
\def\Curvesstack{\mathcal{C}\!\mathit{urves}}
\def\Polarizedstack{\mathcal{P}\!\mathit{olarized}}
\def\Complexesstack{\mathcal{C}\!\mathit{omplexes}}
% \Pic is the operator that assigns to X its picard group, usage \Pic(X)
% \Picardstack_{X/B} denotes the Picard stack of X over B
% \Picardfunctor_{X/B} denotes the Picard functor of X over B
\def\Pic{\mathop{\mathrm{Pic}}\nolimits}
\def\Picardstack{\mathcal{P}\!\mathit{ic}}
\def\Picardfunctor{\mathrm{Pic}}
\def\Deformationcategory{\mathcal{D}\!\mathit{ef}}
\usepackage{cleveref}
\crefname{originalreference}{original reference}{original references}
\Crefname{originalreference}{Original reference}{Original references}
\makeatletter
\newlabel{lemma-finite}{{08YX}{source}{}{sourcepointer-1}{}}
\newlabel{lemma-finite@cref}{{[originalreference][0][]08YX}{[1][0][]source}{}{}{}}
\makeatother
\crefname{originalreference}{original reference}{original references}
\Crefname{originalreference}{Original reference}{Original references}
\makeatletter
\newlabel{lemma-flat}{{00HD}{source}{}{sourcepointer-2}{}}
\newlabel{lemma-flat@cref}{{[originalreference][0][]00HD}{[1][0][]source}{}{}{}}
\makeatother
\crefname{originalreference}{original reference}{original references}
\Crefname{originalreference}{Original reference}{Original references}
\makeatletter
\newlabel{lemma-g1d}{{0CCR}{source}{}{sourcepointer-3}{}}
\newlabel{lemma-g1d@cref}{{[originalreference][0][]0CCR}{[1][0][]source}{}{}{}}
\makeatother
\crefname{originalreference}{original reference}{original references}
\Crefname{originalreference}{Original reference}{Original references}
\makeatletter
\newlabel{lemma-genus-smooth}{{0C1A}{source}{}{sourcepointer-4}{}}
\newlabel{lemma-genus-smooth@cref}{{[originalreference][0][]0C1A}{[1][0][]source}{}{}{}}
\makeatother
\crefname{originalreference}{original reference}{original references}
\Crefname{originalreference}{Original reference}{Original references}
\makeatletter
\newlabel{lemma-grd-inequalities}{{0CCS}{source}{}{sourcepointer-5}{}}
\newlabel{lemma-grd-inequalities@cref}{{[originalreference][0][]0CCS}{[1][0][]source}{}{}{}}
\makeatother
\crefname{originalreference}{original reference}{original references}
\Crefname{originalreference}{Original reference}{Original references}
\makeatletter
\newlabel{lemma-inseparable-linear-system}{{0CD3}{source}{}{sourcepointer-6}{}}
\newlabel{lemma-inseparable-linear-system@cref}{{[originalreference][0][]0CD3}{[1][0][]source}{}{}{}}
\makeatother
\crefname{originalreference}{original reference}{original references}
\Crefname{originalreference}{Original reference}{Original references}
\makeatletter
\newlabel{lemma-linear-series-trivial-existence}{{0CCQ}{source}{}{sourcepointer-7}{}}
\newlabel{lemma-linear-series-trivial-existence@cref}{{[originalreference][0][]0CCQ}{[1][0][]source}{}{}{}}
\makeatother
\crefname{originallemma}{original lemma}{original lemmas}
\Crefname{originallemma}{Original lemma}{Original lemmas}
\makeatletter
\newlabel{morphisms-lemma-characterize-universally-bounded}{{03J5}{source}{}{sourcepointer-8}{}}
\newlabel{morphisms-lemma-characterize-universally-bounded@cref}{{[originallemma][0][]03J5}{[1][0][]source}{}{}{}}
\makeatother
\crefname{originallemma}{original lemma}{original lemmas}
\Crefname{originallemma}{Original lemma}{Original lemmas}
\makeatletter
\newlabel{morphisms-lemma-etale-over-field}{{02GL}{source}{}{sourcepointer-9}{}}
\newlabel{morphisms-lemma-etale-over-field@cref}{{[originallemma][0][]02GL}{[1][0][]source}{}{}{}}
\makeatother
\crefname{originallemma}{original lemma}{original lemmas}
\Crefname{originallemma}{Original lemma}{Original lemmas}
\makeatletter
\newlabel{morphisms-lemma-finite-flat}{{02KB}{source}{}{sourcepointer-10}{}}
\newlabel{morphisms-lemma-finite-flat@cref}{{[originallemma][0][]02KB}{[1][0][]source}{}{}{}}
\makeatother
\crefname{originallemma}{original lemma}{original lemmas}
\Crefname{originallemma}{Original lemma}{Original lemmas}
\makeatletter
\newlabel{morphisms-lemma-finite-locally-free-universally-bounded}{{0CC2}{source}{}{sourcepointer-11}{}}
\newlabel{morphisms-lemma-finite-locally-free-universally-bounded@cref}{{[originallemma][0][]0CC2}{[1][0][]source}{}{}{}}
\makeatother
\title{A smooth proper curve of genus at least 2 has a closed point of separable degree at most 2g - 2}
\author{The Stacks Project Authors}
\date{stacks-project:curves.tex:local-snapshot-20261012; GNU FDL 1.2 or later}
\begin{document}
\maketitle
\noindent\textit{Editorial note.} All delivered mathematical text is the author's own, copied from the pinned Stacks Project chapter file curves.tex; no statement or proof step was rewritten, reordered or removed.\par\medskip
\noindent\textit{Editorial note.} The counted proof is the author's complete argument: the perfect-case argument, the reduction to infinite characteristic p, the Frobenius-twist exclusion of possibility (2), and the closing separability and degree bounds over the etale locus are all retained verbatim.\par\medskip
\noindent\textit{Editorial note.} The label lemma-point-over-separable-extension names the counted statement, and the proof environment immediately following it in the source is delivered verbatim as its proof.\par\medskip
\noindent\textit{Editorial note.} The counted statement is carried by the author's own lemma environment in the chapter source.\par\medskip
\noindent\textit{Allowed standard prerequisites (editorial).}\begin{itemize}
\item The canonical linear series and maps to the projective line. Same chapter, Lemmas lemma-genus-smooth, lemma-linear-series-trivial-existence and lemma-g1d. Gives deg(omega) = 2g - 2 with g global sections, extracts a g\textasciicircum{}1\_\{2g-2\} from the g\textasciicircum{}\{g-1\}\_\{2g-2\}, and converts it into a morphism X -> P\textasciicircum{}1\_k of degree d <= 2g - 2.
\item Finiteness and degree bookkeeping. Same chapter, Lemmas lemma-flat and lemma-finite; Morphisms, Lemmas morphisms-lemma-finite-flat, morphisms-lemma-finite-locally-free-universally-bounded, morphisms-lemma-characterize-universally-bounded and morphisms-lemma-etale-over-field. Makes the morphism finite locally free of degree d, bounds residue-field degrees by d, and gives separability of kappa(u)/kappa(t) over the etale locus.
\item Char p control via the Frobenius twist. Same chapter, Lemmas lemma-grd-inequalities and lemma-inseparable-linear-system. Shows X\textasciicircum{}\{(p)\} carries no g\textasciicircum{}\{g-1\}\_d for d <= g - 1, and excludes possibility (2) of the inseparable-linear-system dichotomy, producing a generically etale map of degree <= 2g - 2.
\end{itemize}\medskip
\ifdefined\hypertarget\hypertarget{sourcepointer-1}{}\fi
\noindent\textit{Source reference.} Original reference 08YX: sources/stacks-curves/pinned-original-source.tex, source lines 323--323. Prerequisite pointer invoked at proof line(s) 2990; it names a standard fact of the same chapter that the counted argument uses, and every step of the counted proof built on it is retained verbatim in the delivered body.\par\medskip
\ifdefined\hypertarget\hypertarget{sourcepointer-2}{}\fi
\noindent\textit{Source reference.} Original reference 00HD: sources/stacks-curves/pinned-original-source.tex, source lines 303--303. Prerequisite pointer invoked at proof line(s) 2989; it names a standard fact of the same chapter that the counted argument uses, and every step of the counted proof built on it is retained verbatim in the delivered body.\par\medskip
\ifdefined\hypertarget\hypertarget{sourcepointer-3}{}\fi
\noindent\textit{Source reference.} Original reference 0CCR: sources/stacks-curves/pinned-original-source.tex, source lines 626--626. Prerequisite pointer invoked at proof line(s) 2984; it names a standard fact of the same chapter that the counted argument uses, and every step of the counted proof built on it is retained verbatim in the delivered body.\par\medskip
\ifdefined\hypertarget\hypertarget{sourcepointer-4}{}\fi
\noindent\textit{Source reference.} Original reference 0C1A: sources/stacks-curves/pinned-original-source.tex, source lines 1760--1760. Prerequisite pointer invoked at proof line(s) 2980; it names a standard fact of the same chapter that the counted argument uses, and every step of the counted proof built on it is retained verbatim in the delivered body.\par\medskip
\ifdefined\hypertarget\hypertarget{sourcepointer-5}{}\fi
\noindent\textit{Source reference.} Original reference 0CCS: sources/stacks-curves/pinned-original-source.tex, source lines 660--660. Prerequisite pointer invoked at proof line(s) 3012; it names a standard fact of the same chapter that the counted argument uses, and every step of the counted proof built on it is retained verbatim in the delivered body.\par\medskip
\ifdefined\hypertarget\hypertarget{sourcepointer-6}{}\fi
\noindent\textit{Source reference.} Original reference 0CD3: sources/stacks-curves/pinned-original-source.tex, source lines 2915--2915. Prerequisite pointer invoked at proof line(s) 3013; it names a standard fact of the same chapter that the counted argument uses, and every step of the counted proof built on it is retained verbatim in the delivered body.\par\medskip
\ifdefined\hypertarget\hypertarget{sourcepointer-7}{}\fi
\noindent\textit{Source reference.} Original reference 0CCQ: sources/stacks-curves/pinned-original-source.tex, source lines 613--613. Prerequisite pointer invoked at proof line(s) 2982; it names a standard fact of the same chapter that the counted argument uses, and every step of the counted proof built on it is retained verbatim in the delivered body.\par\medskip
\ifdefined\hypertarget\hypertarget{sourcepointer-8}{}\fi
\noindent\textit{Source reference.} Original lemma 03J5: sources/stacks-chow/tags.txt, source lines 4627--4627. Prerequisite pointer invoked at proof line(s) 3001; it names a standard fact of a sibling chapter that the counted argument uses, and every step of the counted proof built on it is retained verbatim in the delivered body.\par\medskip
\ifdefined\hypertarget\hypertarget{sourcepointer-9}{}\fi
\noindent\textit{Source reference.} Original lemma 02GL: sources/stacks-chow/tags.txt, source lines 3232--3232. Prerequisite pointer invoked at proof line(s) 3026; it names a standard fact of a sibling chapter that the counted argument uses, and every step of the counted proof built on it is retained verbatim in the delivered body.\par\medskip
\ifdefined\hypertarget\hypertarget{sourcepointer-10}{}\fi
\noindent\textit{Source reference.} Original lemma 02KB: sources/stacks-chow/tags.txt, source lines 3366--3366. Prerequisite pointer invoked at proof line(s) 2992; it names a standard fact of a sibling chapter that the counted argument uses, and every step of the counted proof built on it is retained verbatim in the delivered body.\par\medskip
\ifdefined\hypertarget\hypertarget{sourcepointer-11}{}\fi
\noindent\textit{Source reference.} Original lemma 0CC2: sources/stacks-chow/tags.txt, source lines 15492--15492. Prerequisite pointer invoked at proof line(s) 3000; it names a standard fact of a sibling chapter that the counted argument uses, and every step of the counted proof built on it is retained verbatim in the delivered body.\par\medskip
\begin{lemma}
\label{lemma-point-over-separable-extension}
Let $k$ be a field. Let $X$ be a smooth proper curve over $k$
with $H^0(X, \mathcal{O}_X) = k$ and genus $g \geq 2$.
Then there exists a closed point $x \in X$ with
$\kappa(x)/k$ separable of degree $\leq 2g - 2$.
\end{lemma}

\begin{proof}
Set $\omega = \Omega_{X/k}$. By
Lemma \ref{lemma-genus-smooth} this has degree $2g - 2$
and has $g$ global sections. Thus we have a $\mathfrak g^{g - 1}_{2g - 2}$.
By the trivial Lemma \ref{lemma-linear-series-trivial-existence}
there exists a $\mathfrak g^1_{2g - 2}$
and by Lemma \ref{lemma-g1d} we obtain a morphism
$$
\varphi : X \longrightarrow \mathbf{P}^1_k
$$
of some degree $d \leq 2g - 2$. Since $\varphi$ is flat
(Lemma \ref{lemma-flat}) and finite
(Lemma \ref{lemma-finite})
it is finite locally free of degree $d$
(Morphisms, Lemma \ref{morphisms-lemma-finite-flat}).
Pick any rational point $t \in \mathbf{P}^1_k$
and any point $x \in X$ with $\varphi(x) = t$.
Then
$$
d \geq [\kappa(x) : \kappa(t)] = [\kappa(x) : k]
$$
for example by
Morphisms, Lemmas \ref{morphisms-lemma-finite-locally-free-universally-bounded}
and \ref{morphisms-lemma-characterize-universally-bounded}.
Thus if $k$ is perfect (for example has characteristic zero
or is finite) then the lemma is proved. Thus we reduce to the
case discussed in the next paragraph.

\medskip\noindent
Assume that $k$ is an infinite field of characteristic $p > 0$.
As above we will use that $X$ has a $\mathfrak g^{g - 1}_{2g - 2}$.
The smooth proper curve $X^{(p)}$ has the same genus as $X$.
Hence its genus is $> 0$. We conclude that $X^{(p)}$ does not have a
$\mathfrak g^{g - 1}_d$ for any $d \leq g - 1$ by
Lemma \ref{lemma-grd-inequalities}.
Applying Lemma \ref{lemma-inseparable-linear-system}
to our $\mathfrak g^{g - 1}_{2g - 2}$ (and noting that $2g - 2/p \leq g - 1$)
we conclude that possibility (2) does not occur. Hence we obtain a morphism
$$
\varphi : X \longrightarrow \mathbf{P}^1_k
$$
which is generically \'etale (in the sense of the lemma)
and has degree $\leq 2g - 2$. Let $U \subset X$ be the nonempty
open subscheme where $\varphi$ is \'etale. Then
$\varphi(U) \subset \mathbf{P}^1_k$ is a nonempty Zariski open
and we can pick a $k$-rational point $t \in \varphi(U)$ as $k$ is infinite.
Let $u \in U$ be a point with $\varphi(u) = t$.
Then $\kappa(u)/\kappa(t)$ is separable
(Morphisms, Lemma \ref{morphisms-lemma-etale-over-field}),
$\kappa(t) = k$, and $[\kappa(u) : k] \leq 2g - 2$ as before.
\end{proof}

\end{document}
